\NSPage{073}%
correction cycle of Proposition~\ref{prop:9.6} preserves containment in these prescribed supports; its
estimates require no divisibility by the original cutoff factors.

\section{Oscillatory realization and correction of the residual stress}\label{sec:7}

We fix the background \NSi{NS-F-P073-001}{(u_B,p_B)} of Proposition~\ref{prop:5.5} and the charts, slow cutoffs, and
auxiliary rectangles of Equations~\eqref{eq:6.1}, \eqref{eq:6.9} and \eqref{eq:6.10}. We construct waves whose averaged
quadratic velocity products supply the prescribed radial fluxes of azimuthal and axial
momentum in the leading stress \NSi{NS-F-P073-002}{\mathcal T_0}, with the physical normalization in \eqref{eq:7.30}. Their principal
linearized evolution must also agree with the background shear and viscous diffusion in
\eqref{eq:7.5}. The supports allow us to solve these problems locally and sum the waves without cross
terms between distinct labels, by Lemma~\ref{lem:6.1}. For subsequent corrections, we solve the same
amplitude equation with a prescribed harmonic source using the inverse in Proposition~\ref{prop:7.2}.
A second linear operator, \NSi{NS-F-P073-003}{\mathcal L} in Proposition~\ref{prop:7.6}, prescribes a signed change in the leading
covariance.

The common ingredient is a linear amplitude equation determined by the background
shear and viscosity. We choose phases for which the homogeneous solution \NSi{NS-F-P073-004}{t^{\mathrm h}_{\sigma}} in Lemma~\ref{lem:7.4}
grows and then decays. Its Gaussian bound \eqref{eq:7.16} makes the solution exponentially small
near both ends of its interval. We call this solution a \emph{pulse}. We compute the covariance
of two pulses with different phase gradients and choose their amplitudes to match the
stress. The strict cone condition \eqref{eq:7.1} permits positive squared amplitudes in Proposition~\ref{prop:7.5};
linearizing at those fixed amplitudes permits subsequent stress increments of either sign.
Finally, we take curls of supported vector potentials to enforce exact incompressibility in
Lemma~\ref{lem:7.7}. The additional velocity terms and cutoff errors remain in the covariance estimate
of Corollary~\ref{cor:7.8} and the residual identity \eqref{eq:7.40}.

All calculations use normalized chart representatives until physical factors are displayed.
A band, its representative, its discrete labels, and its integer frequencies are fixed before
differentiation. Constants for derivatives of any fixed order may depend on the fixed profiles
and that order; in applications they may also depend on the finite correction stage. They
remain uniform over bands, labels, and rectangle copies.

\subsection{Phase construction and a frame orthogonal to its gradient}\label{sec:7.1}
We first construct the phase
\NSi{NS-F-P073-005}{\Phi} in \eqref{eq:7.3} that approximately follows transport by the background, controlling the defect in
the leading transport terms as in \eqref{eq:7.9}. For a harmonic \NSi{NS-F-P073-006}{e^{ikm\Phi}}, the leading incompressibility
condition on its amplitude \NSi{NS-F-P073-007}{t_m} is \NSi{NS-F-P073-008}{\nabla_*\Phi\cdot t_m=0}. The frame \NSi{NS-F-P073-009}{B} in \eqref{eq:7.8} for this two-dimensional
plane will compare the undamped principal equation with a diagonal system having one
positive and one negative eigenvalue, as in \eqref{eq:7.10}.

Fix a label \NSi{NS-F-P073-010}{\gamma=(\ell,a,\sigma)} from \eqref{eq:6.8}. On each of its lifted rectangles we use the slow
point \NSi{NS-F-P073-011}{x_*=(R,Z,T)}, transverse coordinate \NSi{NS-F-P073-012}{\xi_g}, and pulse coordinate \NSi{NS-F-P073-013}{v\in[0,L_s]}, where
\NSi{NS-F-P073-014}{L_s=2r_0/c_i\asymp S_*}, as in \eqref{eq:6.11}. Recall \NSi{NS-F-P073-015}{Q=2^{-\ell}}, \NSi{NS-F-P073-016}{\varepsilon=Q^h}, and \NSi{NS-F-P073-017}{S_*=\ell^2} from \eqref{eq:6.1}. The
normalized spatial derivatives are those in \eqref{eq:6.6}; write \NSi{NS-F-P073-018}{\nabla_*f=(D_rf,R^{-1}\partial_\theta f,D_zf)}. A prime
means \NSi{NS-F-P073-019}{\partial_v}, with \NSi{NS-F-P073-020}{x_*} and \NSi{NS-F-P073-021}{\xi_g} fixed. In coefficient estimates, \NSi{NS-F-P073-022}{D^I} denotes ordinary derivatives
in the listed slow and auxiliary coordinates; the individual operators \NSi{NS-F-P073-023}{D_r,D_z} retain their
normalized meanings.

We write the chart base velocity as \NSi{NS-F-P073-024}{be_r+Ve_\theta+Ge_z}, and set \NSi{NS-F-P073-025}{F=V/R} and \NSi{NS-F-P073-026}{g=(RF_R,G_R)}.
Thus \NSi{NS-F-P073-027}{F} is angular velocity, and \NSi{NS-F-P073-028}{g} is the radial shear in the tangential component order \NSi{NS-F-P073-029}{(\theta,z)}.
We identify a tangential two-vector with the vector in \NSi{NS-F-P073-030}{\operatorname{span}(e_\theta,e_z)} having those components.
By \eqref{eq:5.42}, \NSi{NS-F-P073-031}{b=O(\varepsilon)} and the tangential velocity differs from its order-zero value by \NSi{NS-F-P073-032}{O(\varepsilon^2)}, in
\NSPage{074}
every fixed slow derivative. A subscript \NSi{NS-F-P074-001}{0} denotes an order-zero quantity evaluated at the
chosen representative \NSi{NS-F-P074-002}{x^0_{\ell,a}}.

The positive lower bounds in Theorem~\ref{thm:4.6}(iii) allow us to define the following tangential
frame and growth parameters:
\NSFormula{NS-F-P074-003}{display}{none}
\[
N=\frac{g_0}{|g_0|},\qquad K=N^\perp,\qquad
\lambda_0^2=-2F_0N_\theta(2F_0N_\theta+|g_0|)>0,\qquad
c_0=\frac{\lambda_0}{2F_0N_\theta}<0.
\]
We use the quarter turn \NSi{NS-F-P074-004}{K=(-N_z,N_\theta)}, so \NSi{NS-F-P074-005}{N,K} are an orthonormal pair in the tangential
plane, and take \NSi{NS-F-P074-006}{\lambda_0>0}. The definitions of \NSi{NS-F-P074-007}{a,b_s,v_s} in Equations~\eqref{eq:4.11} and \eqref{eq:4.20} give
\NSi{NS-F-P074-008}{g_0=F_0(-a,b_s)} and \NSi{NS-F-P074-009}{\lambda_0^2=2aF_0^2(1-2/v_s)}. In particular, \NSi{NS-F-P074-010}{\lambda_0^2>0} and \NSi{NS-F-P074-011}{c_0<0} follow from the
profile bounds; they are not additional choices.

The phase choice must also retain the stress directions allowed by Lemma~\ref{lem:4.5}. Before
freezing \NSi{NS-F-P074-012}{N,K,c_0}, that pointwise condition is equivalently
\NSFormula{NS-F-P074-013}{display}{7.1}
\begin{equation}
\mathcal T_{0,*}\cdot N<0,\qquad
\left|c_0\frac{\mathcal T_{0,*}\cdot K}{\mathcal T_{0,*}\cdot N}\right|<1,\qquad
\mathcal T_{0,*}=(Q/q)^{A+1/2}\mathcal T_0.
\tag{7.1}\label{eq:7.1}
\end{equation}
For this equivalence, the frame, \NSi{NS-F-P074-014}{c_0}, and target are evaluated at the same point. Indeed,
\NSi{NS-F-P074-015}{c_0^2=(v_s-2)/2}, while the absolute quotient without \NSi{NS-F-P074-016}{c_0} equals \NSi{NS-F-P074-017}{|J_c|/(P_c-v_s)}. At an annulus
boundary the stress vanishes; there the quotient means the one-sided limit obtained from
the smooth unit stress direction in Theorem~\ref{thm:4.6}(iii).

The strict margin on the closed annulus permits us to choose one \NSi{NS-F-P074-018}{u_*>0} such that
\NSi{NS-F-P074-019}{u_*/\sqrt{1+u_*^2}} exceeds the supremum of \NSi{NS-F-P074-020}{|c_0(\mathcal T_{0,*}\cdot K)/(\mathcal T_{0,*}\cdot N)|}, with this interpretation at the
edges. We now freeze \NSi{NS-F-P074-021}{N,K,c_0} at each representative. Uniform continuity of the leading stress
direction preserves the strict inequalities, with a uniform margin, throughout sufficiently
small slow boxes. The frequency and shear bounds, the positive lower bound for \NSi{NS-F-P074-022}{\lambda_0}, and
the frame bounds also hold on fixed small enlargements of the active slow supports. We
take these slow neighborhoods to be enlargements of the mesh boxes by one fixed factor,
relative to \NSi{NS-F-P074-023}{\tau\geq0}, so every point remains within \NSi{NS-F-P074-024}{CS_*^{-3}} of its representative. Derivatives at
\NSi{NS-F-P074-025}{\tau=0} are understood one-sidedly. The uniform cone margin persists on each neighborhood's
intersection with the closed active annulus, using the limiting stress direction at its radial
edges. All these bounds and margins are uniform over bands and labels.

The stress has two components, so we use two phases with different covariance directions
in each slow box. The phase gradient will change along the pulse until viscous damping
exceeds shear amplification. On the rectangle of sign \NSi{NS-F-P074-026}{\sigma\in\{+1,-1\}}, we define
\NSFormula{NS-F-P074-027}{display}{7.2}
\begin{equation}
k=\left\lceil\varepsilon^{-1/2}\right\rceil,\qquad
B_s^2=\frac{2\lambda_0}{\varepsilon k^2(1+u_*^2)^{3/2}},\qquad
s(v)=\sigma\left(\frac{u_*}{2}+\frac{u_*v}{L_s}\right).
\tag{7.2}\label{eq:7.2}
\end{equation}
We take \NSi{NS-F-P074-028}{B_s>0}. The carrier frequency \NSi{NS-F-P074-029}{k} keeps \NSi{NS-F-P074-030}{\varepsilon k^2} of order one, so viscosity remains in the
leading amplitude equation. To choose the angular and axial wave numbers \NSi{NS-F-P074-031}{p,p_z}, we first set
\NSFormula{NS-F-P074-032}{display}{none}
\[
(\widetilde p/R_0,p_z)=B_s\left(K-\frac{\sigma u_*g_0}{L_s|g_0|^2}\right).
\]
Angular periodicity requires \NSi{NS-F-P074-033}{kp} to be an integer. We take \NSi{NS-F-P074-034}{kp} to be a nearest nonzero integer
to \NSi{NS-F-P074-035}{k\widetilde p}, with a fixed rule at ties, and leave \NSi{NS-F-P074-036}{p_z} unchanged. All these wave numbers are constants
for the label. With \NSi{NS-F-P074-037}{x_0=\sigma B_su_*/2}, we define the phase on a lifted rectangle by
\NSFormula{NS-F-P074-038}{display}{7.3}
\begin{equation}
\Phi=p\theta+p_zZ/\varepsilon+x_0R-v(pF+p_zG),
\tag{7.3}\label{eq:7.3}
\end{equation}
\NSFormula{NS-F-P074-039}{display}{7.4}
\begin{equation}
n_\Phi:=\nabla_*\Phi=\bigl(x_0-v(pF_R+p_zG_R),\ p/R,\ p_z-\varepsilon v(pF_Z+p_zG_Z)\bigr).
\tag{7.4}\label{eq:7.4}
\end{equation}
\NSPage{075}
For a source coefficient \NSi{NS-F-P075-001}{f_m\in\mathbb C^3} and a nonzero integer \NSi{NS-F-P075-002}{m}, the principal amplitude problem
is to find \NSi{NS-F-P075-003}{t_m\in\mathbb C^3} and a pressure coefficient \NSi{NS-F-P075-004}{\pi_m\in\mathbb C} satisfying \NSi{NS-F-P075-005}{n_\Phi\cdot t_m=0} and
\NSFormula{NS-F-P075-006}{display}{7.5}
\begin{equation}
t_m'+\mathcal Kt_m+m^2dt_m+ikmn_\Phi\pi_m=-f_m,
\qquad d=\varepsilon k^2|n_\Phi|^2.
\tag{7.5}\label{eq:7.5}
\end{equation}
Here \NSi{NS-F-P075-007}{\mathcal K} is defined below. This equation retains transport in the pulse coordinate, the base
tangential velocity and radial shear, and the viscous term in which both derivatives act on the
exponential. The two factors \NSi{NS-F-P075-008}{2F} in \NSi{NS-F-P075-009}{\mathcal K} include the cylindrical connection terms: radial motion
changes the tangential frame as well as the base swirl. Eliminating the normal pressure term
while differentiating \NSi{NS-F-P075-010}{n_\Phi\cdot t_m=0} gives the projected evolution operator \NSi{NS-F-P075-011}{\mathcal A_\Phi}:
\NSFormula{NS-F-P075-012}{display}{7.6}
\begin{equation}
\mathcal K=
\begin{pmatrix}
0&-2F&0\\
2F+RF_R&0&0\\
G_R&0&0
\end{pmatrix},
\qquad
\mathcal A_\Phi=-\mathcal K+
\frac{n_\Phi\bigl(n_\Phi^T\mathcal K-(n_\Phi')^T\bigr)}{|n_\Phi|^2}.
\tag{7.6}\label{eq:7.6}
\end{equation}
The slow transport, the error in phase transport, and the remaining viscous terms occur in
the residual beyond \eqref{eq:7.5}; Proposition~\ref{prop:9.1} estimates them. For later use, we fix the geometric
coordinates and frame explicitly. Write \NSi{NS-F-P075-013}{n_{\Phi,\mathrm{tan}}=(n_{\Phi,\theta},n_{\Phi,z})}. Since \NSi{NS-F-P075-014}{kp\in\mathbb Z\setminus\{0\}} and \NSi{NS-F-P075-015}{R>0},
the component \NSi{NS-F-P075-016}{n_{\Phi,\theta}=p/R} is nonzero. We may therefore define
\NSFormula{NS-F-P075-017}{display}{7.7}
\begin{equation}
K_a=\frac{n_{\Phi,\mathrm{tan}}}{|n_{\Phi,\mathrm{tan}}|},\qquad
N_a=((K_a)_z,-(K_a)_\theta),\qquad
s_a=\frac{n_{\Phi,r}}{|n_{\Phi,\mathrm{tan}}|}.
\tag{7.7}\label{eq:7.7}
\end{equation}
The corresponding basis of \NSi{NS-F-P075-018}{n_\Phi^\perp} and frame are
\NSFormula{NS-F-P075-019}{display}{7.8}
\begin{equation}
U=[e_r-s_aK_a,N_a],\qquad
B=U
\begin{pmatrix}
1&1\\
c_0\sqrt{1+s^2}&-c_0\sqrt{1+s^2}
\end{pmatrix}.
\tag{7.8}\label{eq:7.8}
\end{equation}
The next lemma compares the actual phase normal and projected operator with their
reference values. In the resulting coordinates, the amplitude equation has two scalar reference
growth rates and a small matrix error.

\begin{lemma}\label{lem:7.1}
For the labels and phases in Equations~\eqref{eq:6.8}, \eqref{eq:7.2} and \eqref{eq:7.3}, there exists a sufficiently
small \NSi{NS-F-P075-020}{q_*>0}, common to all labels and fixed derivative orders, such that every rectangle on \NSi{NS-F-P075-021}{0<q<q_*}
has the following properties. Every \NSi{NS-F-P075-022}{e^{ikm\Phi},\ m\in\mathbb Z\setminus\{0\}}, is single-valued in \NSi{NS-F-P075-023}{\theta}, with nonzero angular
frequency. On the rectangle \NSi{NS-F-P075-024}{0\leq v\leq L_s},
\NSFormula{NS-F-P075-025}{display}{7.9}
\begin{equation}
|n_\Phi-B_s(s(v),K)|+|n_\Phi'|\leq C/S_*,
\qquad
E_{\mathrm{ik}}:=(\mathfrak t_*+bD_r+F\partial_\theta+GD_z)\Phi=O(\varepsilon S_*^C).
\tag{7.9}\label{eq:7.9}
\end{equation}
The estimate for \NSi{NS-F-P075-026}{E_{\mathrm{ik}}} in \eqref{eq:7.9} holds in every fixed coefficient derivative, with a corresponding polynomial
degree. The coefficients extend smoothly to the fixed enlargement. All fixed mixed derivatives in the
slow, transverse, and pulse coordinates of \NSi{NS-F-P075-027}{n_\Phi}, and of the frames used here, have polynomial bounds in
\NSi{NS-F-P075-028}{S_*}.

The real \NSi{NS-F-P075-029}{3\times2} matrix \NSi{NS-F-P075-030}{B} in \eqref{eq:7.8} has columns forming a basis of \NSi{NS-F-P075-031}{n_\Phi^\perp}, and is also viewed as a map
\NSi{NS-F-P075-032}{\mathbb C^2\longrightarrow\{t\in\mathbb C^3:n_\Phi\cdot t=0\}}. It has a uniformly bounded left inverse \NSi{NS-F-P075-033}{B^\ell} and satisfies
\NSFormula{NS-F-P075-034}{display}{7.10}
\begin{equation}
B^\ell(\mathcal A_\Phi B-B')=\operatorname{diag}(\lambda,-\lambda)+E,
\qquad
\lambda(v)=\frac{\lambda_0}{\sqrt{1+s(v)^2}},
\qquad
|E|\leq C/S_*.
\tag{7.10}\label{eq:7.10}
\end{equation}
For the damping coefficient \NSi{NS-F-P075-035}{d} in \eqref{eq:7.5}, put \NSi{NS-F-P075-036}{d_{\mathrm{ref}}=\varepsilon k^2B_s^2(1+s^2)}. Then
\NSFormula{NS-F-P075-037}{display}{7.11}
\begin{equation}
|d-d_{\mathrm{ref}}|\leq C/S_*.
\tag{7.11}\label{eq:7.11}
\end{equation}
\end{lemma}

\begin{proof}
\textbf{Step 1: periodicity and estimates for the phase.}
Because \NSi{NS-F-P075-038}{kp\in\mathbb Z\setminus\{0\}}, the angular
frequency \NSi{NS-F-P075-039}{kmp} is a nonzero integer for every \NSi{NS-F-P075-040}{m\neq0}. This proves the periodicity assertion.
\NSPage{076}
Write \NSi{NS-F-P076-001}{H_\Phi=pF+p_zG}. Since \NSi{NS-F-P076-002}{D_rv=D_zv=0} and the base is independent of the auxiliary
variable,
\NSFormula{NS-F-P076-003}{display}{none}
\[
D_r\Phi=x_0-v(H_\Phi)_R,\qquad
R^{-1}\partial_\theta\Phi=p/R,\qquad
D_z\Phi=p_z-\varepsilon v(H_\Phi)_Z.
\]
These are precisely the three entries of \NSi{NS-F-P076-004}{n_\Phi}. Before rounding, \NSi{NS-F-P076-005}{(p/R_0,p_z)\cdot g_0=-\sigma B_su_*/L_s}.
The slow box has diameter \NSi{NS-F-P076-006}{O(S_*^{-3})}, the base comparison contributes \NSi{NS-F-P076-007}{O(\varepsilon^2)}, and rounding
contributes \NSi{NS-F-P076-008}{O(k^{-1})}. Consequently the error in \NSi{NS-F-P076-009}{pF_R+p_zG_R} is \NSi{NS-F-P076-010}{O(S_*^{-3}+\varepsilon^2+k^{-1})}. Multiplication by \NSi{NS-F-P076-011}{v=O(S_*)}, together with the tangential error \NSi{NS-F-P076-012}{O(S_*^{-1}+\varepsilon S_*)}, proves the first estimate
in \eqref{eq:7.9} after one choice of \NSi{NS-F-P076-013}{q_*}. Differentiation in the pulse coordinate gives \NSi{NS-F-P076-014}{|n_\Phi'|\leq C/S_*}.
Differentiating any fixed number of times introduces only powers of \NSi{NS-F-P076-015}{S_*}; the discrete rounding
is not differentiated. For the phase transport we compute each term:
\NSFormula{NS-F-P076-016}{display}{none}
\[
\mathfrak t_*\Phi=-H_\Phi+\varepsilon v(H_\Phi)_T,
\qquad
F\partial_\theta\Phi+GD_z\Phi=H_\Phi-\varepsilon vG(H_\Phi)_Z.
\]
Consequently
\NSFormula{NS-F-P076-017}{display}{none}
\[
E_{\mathrm{ik}}=\varepsilon v\bigl((H_\Phi)_T-G(H_\Phi)_Z\bigr)+bn_{\Phi,r}.
\]
Every term has a factor \NSi{NS-F-P076-018}{\varepsilon} or \NSi{NS-F-P076-019}{b=O(\varepsilon)}, and the remaining derivatives of fixed order have
polynomial bounds in \NSi{NS-F-P076-020}{S_*}. This proves the differentiated phase-defect estimates as well as
the value estimate. Since \NSi{NS-F-P076-021}{\varepsilon k^2\asymp1} and \NSi{NS-F-P076-022}{B_s} is bounded above and below, \eqref{eq:7.11} follows.

\textbf{Step 2: coordinates on the moving plane.}
It remains to bound the frame \eqref{eq:7.8} and express
the projected evolution in its coordinates. The phase comparison and the lower bound on \NSi{NS-F-P076-023}{B_s}
bound \NSi{NS-F-P076-024}{|n_{\Phi,\mathrm{tan}}|} away from zero. The vector \NSi{NS-F-P076-025}{N_a} in \eqref{eq:7.7} is the perpendicular close to \NSi{NS-F-P076-026}{N}. Let \NSi{NS-F-P076-027}{M}
be the \NSi{NS-F-P076-028}{2\times2} matrix in \eqref{eq:7.8}. Then
\NSFormula{NS-F-P076-029}{display}{none}
\[
B^\ell=M^{-1}
\begin{pmatrix}e_r^T\\N_a^T\end{pmatrix},
\qquad
\begin{pmatrix}e_r^T\\N_a^T\end{pmatrix}U=I_2.
\]
The matrices \NSi{NS-F-P076-030}{U,M,B} and these left inverses are uniformly bounded.

\textbf{Step 3: diagonalization at the reference normal.}
For \NSi{NS-F-P076-031}{n_{\mathrm{ref}}=B_s(s,K)}, write \NSi{NS-F-P076-032}{t=x(e_r-sK)+yN\in n_{\mathrm{ref}}^\perp}, and form \NSi{NS-F-P076-033}{\mathcal K_0} from the reference coefficients. Since \NSi{NS-F-P076-034}{g_0=|g_0|N},
\NSFormula{NS-F-P076-035}{display}{none}
\begin{align*}
(\mathcal K_0t)_r&=2F_0sK_\theta x-2F_0N_\theta y,\\
n_{\mathrm{ref}}\cdot\mathcal K_0t&=2B_sF_0\{K_\theta(1+s^2)x-sN_\theta y\}.
\end{align*}
Let \NSi{NS-F-P076-036}{e=e_r-sK}. The vectors \NSi{NS-F-P076-037}{e,N} are orthogonal, with squared lengths \NSi{NS-F-P076-038}{1+s^2,1}. Directly,
\NSFormula{NS-F-P076-039}{display}{none}
\[
-\mathcal K_0e=(-2F_0sK_\theta,-2F_0e_\theta-g_0),
\qquad
-\mathcal K_0N=(2F_0N_\theta,0),
\]
where each ordered pair gives the radial and tangential parts. Therefore
\NSFormula{NS-F-P076-040}{display}{none}
\[
e\cdot(-\mathcal K_0e)=sK\cdot g_0=0,\qquad
e\cdot(-\mathcal K_0N)=2F_0N_\theta,\qquad
N\cdot(-\mathcal K_0e)=-(2F_0N_\theta+|g_0|).
\]
Dividing the first coordinate by \NSi{NS-F-P076-041}{1+s^2} gives the reference undamped matrix in \NSi{NS-F-P076-042}{(x,y)}:
\NSFormula{NS-F-P076-043}{display}{none}
\[
\begin{pmatrix}
0&2F_0N_\theta/(1+s^2)\\
-(2F_0N_\theta+|g_0|)&0
\end{pmatrix}.
\]
Its eigenvectors are the columns of \NSi{NS-F-P076-044}{M}, with eigenvalues \NSi{NS-F-P076-045}{\lambda,-\lambda}. To compare this matrix with
the actual evolution, first use the value estimates for the base coefficients and \NSi{NS-F-P076-046}{n_\Phi-n_{\mathrm{ref}}}. They
change the algebraic projection \NSi{NS-F-P076-047}{-\mathcal K+n_\Phi n_\Phi^T\mathcal K/|n_\Phi|^2} by \NSi{NS-F-P076-048}{O(S_*^{-1})}. Also \NSi{NS-F-P076-049}{K_a-K} and \NSi{NS-F-P076-050}{s_a-s} are
\NSi{NS-F-P076-051}{O(S_*^{-1})}, so the actual frame and its left inverse differ from their reference values by that
order. Expressing the algebraic projection in these frames therefore changes the reference
matrix by \NSi{NS-F-P076-052}{O(S_*^{-1})}. The moving-normal term \NSi{NS-F-P076-053}{-n_\Phi(n_\Phi')^T/|n_\Phi|^2} has the same bound, since
\NSi{NS-F-P076-054}{|n_\Phi'|=O(S_*^{-1})} and \NSi{NS-F-P076-055}{|n_\Phi|} is bounded below. Also \NSi{NS-F-P076-056}{s'=O(S_*^{-1})}, so the definitions of \NSi{NS-F-P076-057}{U,M} give
\NSPage{077}
\NSi{NS-F-P077-001}{U',M'=O(S_*^{-1})}. Hence \NSi{NS-F-P077-002}{B'=U'M+UM'=O(S_*^{-1})}, and multiplication by the bounded
\NSi{NS-F-P077-003}{B^\ell} gives \NSi{NS-F-P077-004}{-B^\ell B'=O(S_*^{-1})}. These contributions give the error \NSi{NS-F-P077-005}{E} in \eqref{eq:7.10}.

We now specify why this choice of \NSi{NS-F-P077-006}{q_*} works for every fixed derivative order. The value
comparisons in Step 1 require only finitely many smallness conditions. For example, we can
require
\NSFormula{NS-F-P077-007}{display}{none}
\[
S_*^2(\varepsilon+\varepsilon^2+k^{-1})\leq1
\]
to bound the errors multiplied by \NSi{NS-F-P077-008}{v=O(S_*)} by \NSi{NS-F-P077-009}{C/S_*}. Since \NSi{NS-F-P077-010}{\varepsilon=2^{-h\ell}}, \NSi{NS-F-P077-011}{k^{-1}\leq\varepsilon^{1/2}}, and
\NSi{NS-F-P077-012}{S_*=\ell^2}, this holds for all sufficiently large band indices. The bound \NSi{NS-F-P077-013}{1\leq\varepsilon k^2\leq4} for \NSi{NS-F-P077-014}{\varepsilon\leq1}
gives fixed positive upper and lower bounds for \NSi{NS-F-P077-015}{B_s}. We make the phase-normal error smaller
than half this lower bound. The estimate \NSi{NS-F-P077-016}{|n_{\Phi,\mathrm{tan}}|\geq B_s-C/S_*} then keeps the tangential
phase normal uniformly nonzero; the same lower bound applies to \NSi{NS-F-P077-017}{|n_\Phi|}. The formula
\NSi{NS-F-P077-018}{\det M=-2c_0\sqrt{1+s^2}} and the fixed lower bound for \NSi{NS-F-P077-019}{|c_0|} likewise bound \NSi{NS-F-P077-020}{|\det M|} away from
zero. Thus one lower bound on the band index, or equivalently one upper bound \NSi{NS-F-P077-021}{q_*} on \NSi{NS-F-P077-022}{q},
makes every denominator in the frame and its left inverse uniformly nonzero.

With that domain fixed, any prescribed number of coefficient derivatives of the base
coefficients and the phase normal \NSi{NS-F-P077-023}{n_\Phi} has a bound \NSi{NS-F-P077-024}{C_I S_*^{b_I}}. Applying the product and quotient
rules to the displayed frame formulas gives bounds of the same form for \NSi{NS-F-P077-025}{U,M,B,B^\ell} and
their pulse derivatives, using the lower bounds just established. Differentiating the phase
defect retains its factor \NSi{NS-F-P077-026}{\varepsilon}: the band is fixed, and every fixed derivative of \NSi{NS-F-P077-027}{b} is \NSi{NS-F-P077-028}{O(\varepsilon)}. These
arguments require bounds on the higher derivatives, with constants and polynomial degrees
allowed to depend on \NSi{NS-F-P077-029}{I}; they impose no further smallness condition on those derivatives.
Consequently they use the same \NSi{NS-F-P077-030}{q_*}.
\end{proof}

\subsection{Solving the pulse equation and controlling derivatives}\label{sec:7.2}
The frame \NSi{NS-F-P077-031}{B} from Lemma~\ref{lem:7.1}
reduces \eqref{eq:7.5} to an ordinary differential equation along each pulse. In this subsection we
will first solve it with zero initial amplitude and a prescribed source \NSi{NS-F-P077-032}{f_m}, obtaining the
linear inverse used to remove later harmonics in Proposition~\ref{prop:7.2}. We will then solve the
zero-source equation with a specified nonzero initial amplitude to produce the pulses for
the leading covariance in Lemma~\ref{lem:7.4}. Localizing both solutions requires decay at both ends
of the interval. Subtracting the reference damping from the growing eigenvalue defines the
common envelope
\NSFormula{NS-F-P077-033}{display}{7.12}
\begin{equation}
P(v)=\exp\left(\int_{L_s/2}^{v}(\lambda(w)-d_{\mathrm{ref}}(w))\,dw\right).
\tag{7.12}\label{eq:7.12}
\end{equation}
We also write \NSi{NS-F-P077-034}{P_v=P(v)}, as in the coefficient class \eqref{eq:6.29}. The estimates \eqref{eq:7.14} below show
that a source bounded by this envelope has a response with the same envelope, at the cost
of polynomial factors in \NSi{NS-F-P077-035}{S_*}. Multiplying a solution by \NSi{NS-F-P077-036}{\psi} will therefore change its principal
equation only where the solution and source are exponentially small, as shown after \eqref{eq:7.40}.

\begin{proposition}\label{prop:7.2}
Fix \NSi{NS-F-P077-037}{\alpha\in\mathbb R}, a finite set of harmonics \NSi{NS-F-P077-038}{0<|m|\leq M}, and a label at level \NSi{NS-F-P077-039}{i} on a
common torus \NSi{NS-F-P077-040}{H=Y_{i_0}}, with \NSi{NS-F-P077-041}{i-i_0} bounded as in Lemma~\ref{lem:6.2}. For each harmonic, let \NSi{NS-F-P077-042}{f_m(x_*,H)\in\mathbb C^3}
be a coefficient descending to \NSi{NS-F-P077-043}{H}, smooth on the fixed enlargement of that label's rectangles. Assume
that it has a fixed closed slow support contained in the label's slow support, a compact transverse
support contained in \NSi{NS-F-P077-044}{\supp\chi_g}, and support in the closed active shell \NSi{NS-F-P077-045}{X_a\leq X\leq X_b}, with smooth zero
extension across all these support boundaries. Assume also that on \NSi{NS-F-P077-046}{X_a<X<X_b} and \NSi{NS-F-P077-047}{0\leq v\leq L_s},
for every fixed multi-index \NSi{NS-F-P077-048}{I} of derivatives in \NSi{NS-F-P077-049}{(R,Z,T,H)},
\NSFormula{NS-F-P077-050}{display}{none}
\[
|D^If_m|\leq C_I\varepsilon^\alpha S_*^{b_I}\sqrt{\zeta}\,\delta^{-a_I}P(v).
\]
\NSPage{078}
On the domain \NSi{NS-F-P078-001}{0<q<q_*} fixed in Lemma~\ref{lem:7.1}, there is a unique amplitude \NSi{NS-F-P078-002}{t_m} satisfying \NSi{NS-F-P078-003}{n_\Phi\cdot t_m=0},
with zero initial data, and a pressure coefficient \NSi{NS-F-P078-004}{\pi_m} such that \eqref{eq:7.5} holds. They are given by
\NSFormula{NS-F-P078-005}{display}{7.13}
\begin{equation}
\begin{aligned}
t_m'&=\mathcal A_\Phi t_m-m^2dt_m-\operatorname{proj}_{n_\Phi^\perp}f_m,
&\qquad t_m(0)&=0,\\
\pi_m&=\frac{i}{km}\,
\frac{n_\Phi\cdot\mathcal Kt_m-n_\Phi'\cdot t_m+n_\Phi\cdot f_m}{|n_\Phi|^2}.
\end{aligned}
\tag{7.13}\label{eq:7.13}
\end{equation}
Here \NSi{NS-F-P078-006}{\operatorname{proj}_{n_\Phi^\perp}f=f-n_\Phi(n_\Phi\cdot f)/|n_\Phi|^2} is orthogonal projection. The source is evaluated along the
path \NSi{NS-F-P078-007}{H_w} in Equation~\eqref{eq:6.21}, with \NSi{NS-F-P078-008}{x_*=(R,Z,T)} and the transverse coordinate fixed. The coefficients
descend to the same common torus and preserve the specified slow, transverse, and shell supports,
including smooth zero extension there. The solution is defined on the entire closed pulse interval
\NSi{NS-F-P078-009}{0\leq v\leq L_s}. On \NSi{NS-F-P078-010}{X_a<X<X_b}, for every fixed \NSi{NS-F-P078-011}{I},
\NSFormula{NS-F-P078-012}{display}{7.14}
\begin{equation}
|D^It_m|\leq C_I'\varepsilon^\alpha S_*^{b_I'}\sqrt{\zeta}\,\delta^{-a_I'}P(v),
\tag{7.14}\label{eq:7.14}
\end{equation}
\NSFormula{NS-F-P078-013}{display}{7.15}
\begin{equation}
|D^I\pi_m|\leq C_I'\varepsilon^{\alpha+1/2}S_*^{b_I'}\sqrt{\zeta}\,\delta^{-a_I'}P(v).
\tag{7.15}\label{eq:7.15}
\end{equation}
Constants and degrees may depend on \NSi{NS-F-P078-014}{M,I}, and the prescribed source bounds. The threshold \NSi{NS-F-P078-015}{q_*}
does not depend on \NSi{NS-F-P078-016}{M,\alpha}, or the correction stage. If the source has uniform one-sided derivatives in
the slow variables at \NSi{NS-F-P078-017}{\tau=0} on compact sets with \NSi{NS-F-P078-018}{q>0}, the solution and pressure have the same
one-sided derivatives.
\end{proposition}

Only descent to the common torus \NSi{NS-F-P078-019}{H} is required in Proposition~\ref{prop:7.2}: the source may take
different values at distinct preimages of a band-torus point. No orthogonality to \NSi{NS-F-P078-020}{n_\Phi} or
divisibility by either rectangle cutoff is assumed. The support conclusion concerns the slow,
transverse, and shell variables; no temporal zero extension is asserted before multiplication
by \NSi{NS-F-P078-021}{\psi}.

\begin{proof}
The moving frame gives a two-dimensional equation. We first bound its forward
propagator by \NSi{NS-F-P078-022}{P(v)/P(w)}, then use the resulting Duhamel formula for values and derivatives.
Finally, we check support, descent, and the normal pressure equation.

\textbf{Step 1: the envelope and forward propagator.}
The scalar reference growth rate is
\NSFormula{NS-F-P078-023}{display}{none}
\[
a_{\mathrm{net}}(y)=\frac{\lambda_0}{\sqrt{1+y^2}}
-\frac{\lambda_0(1+y^2)}{(1+u_*^2)^{3/2}},
\qquad y=|s(v)|.
\]
It vanishes at \NSi{NS-F-P078-024}{y=u_*}, and its derivative in the pulse coordinate is
\NSFormula{NS-F-P078-025}{display}{none}
\[
\frac{d}{dv}a_{\mathrm{net}}(|s(v)|)
=-\frac{u_*\lambda_0|s(v)|}{L_s}
\left((1+s(v)^2)^{-3/2}+2(1+u_*^2)^{-3/2}\right).
\]
Because \NSi{NS-F-P078-026}{u_*/2\leq|s(v)|\leq3u_*/2}, this derivative lies between \NSi{NS-F-P078-027}{-C/L_s} and \NSi{NS-F-P078-028}{-c/L_s}. By \eqref{eq:7.12},
\NSi{NS-F-P078-029}{(\log P)'=a_{\mathrm{net}}(|s(v)|)}, and both \NSi{NS-F-P078-030}{\log P} and its first derivative vanish at \NSi{NS-F-P078-031}{v=L_s/2}. Integrating
these derivative bounds twice from the midpoint gives
\NSFormula{NS-F-P078-032}{display}{7.16}
\begin{equation}
e^{-C(v-L_s/2)^2/L_s}\leq P(v)\leq e^{-c(v-L_s/2)^2/L_s}\leq1.
\tag{7.16}\label{eq:7.16}
\end{equation}
The envelope therefore decays on both sides of the midpoint. To show that the sourced
amplitude \NSi{NS-F-P078-033}{t_m} has the same decay, we write \NSi{NS-F-P078-034}{t_m=Bz_m} using the actual moving frame. The
coordinate vector \NSi{NS-F-P078-035}{z_m} satisfies
\NSFormula{NS-F-P078-036}{display}{7.17}
\begin{equation}
z_m'=A_mz_m+g_m,
\qquad A_m=\operatorname{diag}(\lambda,-\lambda)+E-m^2dI_2,
\qquad g_m=-B^\ell\operatorname{proj}_{n_\Phi^\perp}f_m.
\tag{7.17}\label{eq:7.17}
\end{equation}
